% エントリ 9(2026-11-03): 独立な検証と訂正
\section{[E9, 2026-11-03] 独立な検証と訂正}\label{sec:e09}

E1--E8 を三つの観点から独立に検証した(ユーザー指示「研究ノートとコードを丁寧にチェック」)。
\begin{itemize}
\item \textbf{理論}: 導出・規約・近似の妥当範囲を原論文と照合した。
\item \textbf{コード}: 式と実装を一行ずつ比べ、全スクリプトを再実行して本文の数値と照合した。
\item \textbf{観測}: 数値と主張を原論文の本文と照合した。
\end{itemize}
検証者には結論を伝えず、ファイルだけを渡した。過去のエントリは書き換えず、訂正をここにまとめる。計算は \texttt{calc/e09\_review.py}、出力は \texttt{calc/out/e09.json}。

\subsection{確かめられたこと}

\begin{itemize}
\item \textbf{規約と式}: \arxiv{1908.10370} の式 2.32 は振幅の率で、占有数は $2\Gamma$ で育つこと。$\tilde r_+$、$\Mpl=1/\sqrt G$、$\varepsilon=N/GM^2$。BGLS の準平衡(式 55a,b)と、回転を失う速さが $f^2$ に比例すること。E5 の運動論の式(BGLS の式 34--36 と、多重度・符号・係数まで一致)。E4 の Dolan の係数 $c_0$--$c_4$ と角度の固有値。
\item \textbf{保存則}: 準平衡では、超放射で取り出した角運動量が $m=3$ の放出で外へ出ていく量と等しい。エネルギーも $3\omega_{211}=2(2\omega_{211}-\omega_{322})+(2\omega_{322}-\omega_{211})$ で閉じる。
\item \textbf{年齢の比}: 準平衡の中で、年齢の比が $f$ によらないこと。$R$ の値も独立に再現された。
\item \textbf{コード}: 全スクリプトを再実行した出力は commit 済みのものと一致した。本文の数値は、下の訂正を除いてすべて再現した。
\item \textbf{文献の数値}: \arxiv{2607.01317}・\arxiv{2609.02678} の結果。GW241011 の主星($m_1=19.6^{+3.6}_{-2.5}\Msun$、$\chi_1=0.78\pm0.09$、0.69 は 95\% の下限)。\arxiv{2605.27226} の速く回る部分集団。Cyg X-1 の質量と fiducial のスピン。
\item \textbf{$C_\mathrm{BH}$ の表}(E7): \arxiv{2412.03655} の Fig.~3 をベクトル図から読み直した値と 8\% 以内で一致した。
\end{itemize}

\subsection{訂正 1(最重要): 窓は仮定に強く依存する}

E8 の窓は、三つの楽観的な仮定を重ねたものだった。
\begin{enumerate}
\item \textbf{GRS 1915+105 の年齢とスピン}: E8 は \arxiv{2412.03655} にならって、年齢を Salpeter 時間 $4.5\times10^7\yr$ とした。これは、降着でスピンが戻りうることを前提にしている。
  \begin{itemize}
  \item この連星の伴星は $\sim0.5\Msun$ と軽い。スピンを臨界値 0.26 から 0.95 まで戻すには $\sim5\Msun$ の降着が要る(Bardeen の薄い円盤の降着で $M_f/M_i=1.63$。0.87 から 0.92 でも 1.074 倍、約 $0.9\Msun$)ので、現実的ではない。
  \item 制限を出すには短い年齢が控えめだが、窓が残ると主張するには楽観的である。
  \item スピンの下限は、\arxiv{2412.03655} の 0.95 ではなく、元の測定(Reid+ 2014、\arxiv{1409.2453})の系統誤差込みの 3$\sigma$ で 0.92 である。
  \end{itemize}
\item \textbf{GW241011 の生まれたときのスピン}: E3--E8 は観測の中央値 0.78 を初期値とした。2 世代目の予言は $\approx0.69$(0.53--0.72、E8)なので、0.72 で生まれたなら許される回転の減少は 0.03 しかない。
\item \textbf{2 世代目の年齢}: E8 の表の「球状星団で典型 10--100 Myr」は、\arxiv{2103.05016} の Fig.~2 から読んだ。しかしその量は、前の世代の形成を含む累積の時間だった(同論文の式 21)。重い 2 世代目の合体までの遅れの最小値は、$\sim5\times10^7\yr$($50+50\Msun$)から $3\times10^8\yr$($10+10\Msun$)である(Baibhav+ \arxiv{2004.00650} の Fig.~2 からの読み取り)。E8 の窓が要る $T_2\lesssim3\times10^7\yr$ より長い。
\end{enumerate}

そこで、次の二通りの仮定で窓を描き直した(図~\ref{fig:e09})。
\begin{itemize}
\item \textbf{楽観的}(E8 と同じ): GRS は 0.995 から 0.95 まで $4.5\times10^7\yr$ で落ちない。GW241011 は 0.78 で生まれ、$T_2=3\times10^7\yr$。
\item \textbf{現実的}: GRS は BH の年齢 $3\times10^9\yr$ のあいだ 0.87 を保つ(下限 0.92 から、軽い伴星からの降着で戻せる $\sim0.05$ を差し引く)。GW241011 は 0.72 で生まれ、$T_2=5\times10^7\yr$。
\end{itemize}

\begin{figure}[H]
\centering
% E9: 楽観的な仮定と現実的な仮定での窓(gamma_BH は主要項)。データは calc/e09_review.py
\begin{tikzpicture}
\begin{groupplot}[group style={group size=2 by 1, horizontal sep=14mm},
  width=0.48\linewidth, height=0.46\linewidth, note, xmode=log, ymode=log,
  xmin=0.4e-12, xmax=1.2e-12, ymin=3e10, ymax=1e13,
  xtick={0.4e-12,0.6e-12,0.8e-12,1e-12}, xticklabels={0.4,0.6,0.8,1.0},
  xlabel={$\mu\ [10^{-12}\,\mathrm{eV}]$}]
\nextgroupplot[ylabel={$f\ [\mathrm{GeV}]$}, title={\small (a) 楽観的な仮定(E8)},
  legend style={font=\scriptsize, at={(0.02,0.02)}, anchor=south west, fill=white, fill opacity=0.85, text opacity=1, draw=none}, legend cell align=left]
  \fill[cB!10] (axis cs:0.60e-12,3e10) rectangle (axis cs:0.89e-12,1e13);
  \addplot[ink, line width=1.0pt] table[x=mu, y=GRS] {figs/data/e09_window_opt.dat}; \addlegendentry{GRS 1915+105(上を排除)}
  \addplot[cA, line width=1.0pt] table[x=mu, y=GW241011] {figs/data/e09_window_opt.dat}; \addlegendentry{GW241011(上を排除)}
  \addplot[cC!60!black, dashed, line width=0.9pt] table[x=mu, y=M15_d18] {figs/data/e09_window_opt.dat}; \addlegendentry{1 世代目が失う(下を排除): $15\Msun$}
  \addplot[cC, dashed, line width=0.9pt] table[x=mu, y=M20_d18] {figs/data/e09_window_opt.dat}; \addlegendentry{\hphantom{1 世代目が失う(下を排除): }$20\Msun$}
  \addplot[cC!45, dashed, line width=0.9pt] table[x=mu, y=M25_d18] {figs/data/e09_window_opt.dat}; \addlegendentry{\hphantom{1 世代目が失う(下を排除): }$25\Msun$}
\nextgroupplot[title={\small (b) 現実的な仮定}]
  \fill[cB!10] (axis cs:0.60e-12,3e10) rectangle (axis cs:0.89e-12,1e13);
  \addplot[ink, line width=1.0pt] table[x=mu, y=GRS] {figs/data/e09_window_real.dat};
  \addplot[cA, line width=1.0pt] table[x=mu, y=GW241011] {figs/data/e09_window_real.dat};
  \addplot[cC!60!black, dashed, line width=0.9pt] table[x=mu, y=M15_d18] {figs/data/e09_window_real.dat};
  \addplot[cC, dashed, line width=0.9pt] table[x=mu, y=M20_d18] {figs/data/e09_window_real.dat};
  \addplot[cC!45, dashed, line width=0.9pt] table[x=mu, y=M25_d18] {figs/data/e09_window_real.dat};
\end{groupplot}
\end{tikzpicture}

\caption{仮定の二つの組での窓($\gamma_\mathrm{BH}$ は主要項)。実線の下で GRS 1915+105 と GW241011 が回転を保ち、破線の上で 1 世代目($\chi_0=0.7$)が 6 Gyr のうちに 211 の臨界 $+0.18$ まで回転を失う。橙の帯は \arxiv{2609.02678} の質量。}
\label{fig:e09}
\end{figure}

\begin{result}
$\mu=0.71\times10^{-12}\eV$ では次のようになる(主要項の $\gamma_\mathrm{BH}$。括弧は E7 の補正を入れた値)。
\begin{itemize}
\item \textbf{楽観的}: 上限は $f\le2.3\times10^{12}\GeV$(GRS、$2.9\times10^{12}$)と $2.4\times10^{12}\GeV$(GW241011、$3.9\times10^{12}$)。1 世代目が回転を失う下限($\delta=0.18$)は、$10$、$15$、$20\Msun$ で $1.7$、$0.56$、$0.22\times10^{12}\GeV$。$10\Msun$ 以上の 1 世代目が窓に入る。
\item \textbf{現実的}: 上限は $f\le5.1\times10^{11}\GeV$(GRS、$6.4\times10^{11}$)と $1.1\times10^{12}\GeV$(GW241011、$1.8\times10^{12}$)。下限は $15\Msun$ で $5.6\times10^{11}$($7.6\times10^{11}$)、$20\Msun$ で $2.2$--$4.6\times10^{11}$、$25\Msun$ で $0.7$--$2.4\times10^{11}\GeV$($\delta=0.18$--$0.05$)。窓が残るのは、$f\approx(2$--$5)\times10^{11}\GeV$ で、回転を失うのが $\sim20\Msun$ 以上の 1 世代目だけの場合である。
\item \textbf{仮定を一つずつ変えた場合}(GRS の上限、主要項):
  \begin{itemize}
  \item $\tilde a>0.92$・$\tau_\mathrm{Sal}$: $3.1\times10^{12}\GeV$
  \item $\tilde a>0.95$・3 Gyr: $2.8\times10^{11}\GeV$
  \item 0.87・5 Gyr: $4.0\times10^{11}\GeV$
  \end{itemize}
  GW241011 が 0.70 で生まれ、$T_2=5\times10^7\yr$ なら $6.6\times10^{11}\GeV$。
\item \textbf{Cyg X-1 の fiducial のスピン}: 相対論的な率で $3.5\times10^{11}\GeV$(E8 の $2.9\times10^{11}$ は解析式の率による)。最初のスピンを Thorne の上限 0.998 にして 0.997 まで、としても $3.5\times10^{11}\GeV$。現実的な GRS の上限と同じ程度になる。
\end{itemize}
\end{result}

\paragraph{意味}
E8 の「$f\approx(1$--$2.5)\times10^{12}\GeV$ の細い帯」は、楽観的な仮定をすべて重ねたときだけ成り立つ。現実的な仮定では、窓は $f\lesssim5\times10^{11}\GeV$ に押し下げられ、回転を失うのは $\sim20\Msun$ 以上の 1 世代目だけになる。集団解析の問いは次の形になる。
\begin{itemize}
\item GWTC-5 の積み上がりが $\sim20\Msun$ 以上に限られることを許すか。
\item 許さないなら、\arxiv{2609.02678} の質量は自己相互作用があっても排除される。
\end{itemize}

\subsection{訂正 2: E6 の論証}

E6 の「$\chi\ge\chi_c^{(211)}$ なら $m\ge1$ のモードは吸収されない」は、一般には正しくない。$\chi_c^{(211)}$ と、$\Omega_H=\mu$ となるスピン $\chi_\mu$ のあいだでは、$\omega_{211}<\omega<\mu$ の $m=1$ のモード(311 などの $n\ge3$ の準位や、強制振動)は吸収され、式 \eqref{eq:e06_dchi} の括弧が $1-2\chi M\omega>0$ なのでスピンを上げる。結論(吸収でスピンは戻らない)は、次の論証で保たれる。
\begin{itemize}
\item \textbf{上限 1}: $m=1$ のモードの吸収は $\Omega_H$ を上げ、$\Omega_H=\omega<\mu$ で止まる(そこでそのモードは超放射する側になる)。したがって吸収で上がるのは $\chi_\mu$ までである。$\chi_\mu-\chi_c^{(211)}$ は、$\mu=0.71\times10^{-12}\eV$ で GW241011 $5.0\times10^{-4}$、GRS $1.4\times10^{-4}$、Cyg X-1 $6.2\times10^{-4}$、$10\Msun$ $7.3\times10^{-5}$ である。
\item \textbf{上限 2}(外からのトルクがない場合): $\chi=J/M^2$ は $J$ について増加し、既約質量 $M_\mathrm{irr}$ について減少する。面積定理で $M_\mathrm{irr}$ は減らず、雲の正の角運動量はすべて BH から来たので $J_\mathrm{BH}\le J_0$。よって $\chi\le\chi_0$。
\item \textbf{準平衡}: 自己相互作用の準平衡では、取り出した角運動量は $m=3$ のスカラー波として外へ出ていくので、戻らない。
\item \textbf{例外}: ガスの降着と、伴星の軌道からのトルク(外からの角運動量)。
\end{itemize}
DESIGN.md D2 の「減衰する準位に移った雲は $m\le$ 元の $m$ の角運動量を返すだけ」も同じ理由で不正確なので、直す。

\subsection{訂正 3: BH の質量の変化}

E3・E5 は、雲が小さいことを理由に BH の質量の変化を落とした。しかし準平衡では、超放射で取り出したエネルギーの 2/3 を BH が $211\times211\to322\times\mathrm{BH}$ で吸い戻す。そのため BH の質量は、取り出した量の 1/3 だけ減る。このとき、
\[
\dot\chi=\frac{\dot J}{M^2}\Bigl(1-\frac{2\chi\alpha}{3}\Bigr)
\]
となり、雲の大きさによらない。回転を失う時間は GW241011 で 5.4\%、GRS で 4.5\%、1 世代目($10$、$15\Msun$)で 1.8\%、3.2\% 長くなる。$f$ の上限・下限は 1--3\% 動く。\texttt{si.MASS\_CORR} で入れられるようにした。本計算では入れる。

\subsection{訂正 4: 数値と記述}

\begin{itemize}
\item \textbf{連分数の打ち切り}(E4): 「$\alpha\ge0.03$ では 4000 項と 16000 項で 5 桁以上一致」は、高いスピンで成り立たない。$\alpha=0.031$、$\chi=0.995$ で 4.7\% ずれた。項数を $\max(4000,500/\alpha)$ に増やし、表を作り直した。$N$ と $2N$ の差は 2.3$\times10^{-4}$ 以下。E3--E8 の引用した数値は変わらない(変化は $10^{-6}$ 以下)。
\item \textbf{433 の表}: 使える範囲を $0.30\le\alpha\le0.40$ に狭めた(端で多倍長との差が 18\%、欠けた点があった)。引用した数値には効かない。
\item \textbf{E4 の記述}:
  \begin{itemize}
  \item 「解析式は $\alpha\lesssim0.08$ で 10\% 以内」は、$\chi=0.5$ では $\alpha\lesssim0.076$ まで。
  \item \arxiv{1908.10370} の「$\alpha\lesssim0.2$ で優れた近似」は、エネルギー(式 2.28)についての記述で、率(式 2.32)についてではない(E2 の記述も同じ)。
  \item Dolan の式の記号 $\chi=(\mu^2-2\omega^2)/q$ はスピンと紛らわしい。
  \item 同じ時間を E2 では 301 年、E4 では 302 年と書いた(時間の格子の違い)。
  \end{itemize}
\item \textbf{E7}: 「補正は比を最大 2 倍動かす」は、2.2 倍($M_1=8\Msun$)から 2.7 倍($\mu=0.89\times10^{-12}\eV$)まで。
\item \textbf{E8}:
  \begin{itemize}
  \item 「GRS の値は 10\% で合う」は、比べる \arxiv{2412.03655} の計算による。Fig.~1(二つの BH をまとめたもの)と Fig.~14 とは 10--15\%、Fig.~16 右(GRS だけ、$n\le5$、$1.5\times10^{12}$)とは 65\%。
  \item JSON の要約は、最も近い格子点($0.69\times10^{-12}$)の値だった。ちょうどその $\mu$ で計算するよう直した。
  \item Cyg X-1 の年齢 4.8--7.6 Myr は、$14.8\Msun$ の BH と $19.2\Msun$ の伴星で出された値(\arxiv{1107.5585})。重い伴星なら若くなり、制限は弱まる。
  \item Cyg X-1 の控えめなスピンは、Zdziarski+ \arxiv{2402.12325} の版で $0.92^{+0.07}_{-0.05}$ から $0.87$--$0.90$ まで変わる。
  \item 孤立連星の年齢の「中央値 2.8 Gyr」は、形成の経路によらない GWTC-2 の値(\arxiv{2105.06491})。
  \item 「$\gtrsim6$ Gyr」は、すべての星団を 12 Gyr、$z<0.5$ とした場合(\arxiv{1911.00018} の Fig.~15)。
  \item ずれ機構で暗黒物質のすべてになる $f$ は、$\theta_i\sim1$ で $f\theta_i\approx(3$--$7)\times10^{14}\GeV$($3H=\mu$ と $H=\mu$ で振動が始まる場合。\texttt{e09.json})。E8 の $2\times10^{13}\GeV$ は、$\theta_i=\pi-10^{-6}$ まで調整した極端な値だった。$f=2\times10^{12}\GeV$、$\theta_i=1$ では暗黒物質の $\sim4\times10^{-5}$ になる。
  \end{itemize}
\item \textbf{E1}:
  \begin{itemize}
  \item 「分かれる原因は幅と一つの事象」は単純化しすぎた。\arxiv{2609.02678} のデルタ関数の計算は、幅を 0 にしただけでなく、次の点も変えている(同論文の End Matter III)。
    \begin{itemize}
    \item 基準のスピン模型を Beta にした。
    \item $T_\mathrm{age}=10\,\mathrm{Myr}$ に固定した。
    \item 質量模型を単一のべき乗と一様な $q$ にした。
    \item KDE でスピンの大きさだけを使った。
    \end{itemize}
    \arxiv{2607.01317} とは、選択効果の有無、年齢の分布、率の扱いも違う。
  \item GW241011 を 2 世代目とする解釈の典拠は、\review\ T04-28 ではなく \arxiv{2510.26931} の \S4.2(T04-29)。
  \end{itemize}
\item \textbf{E2}: 複数の準位の半解析は、322 が $\chi$ を 211 の臨界より下げたときの 211 の吸い戻しを落としている。最終のスピンは変わらないが、322 の飽和の時刻が $e$ 倍 1 回分ほどずれる。
\item \textbf{未確認}: GW190517($M_1\approx39\Msun$、$\chi_1\approx0.9$)は、自己相互作用の制限(\arxiv{2507.20979})が $\alpha<0.2$ の選び方で $5\times10^{-13}\eV$ までしか出していない。この $\mu$ で効くかを、このノートの模型で確かめる。
\end{itemize}
